\section{Introduction}
\label{sec:introduction}

Topologically close-packed (TCP) phases are ubiquitous in refractory transition-metal alloys and strongly influence the long-term stability, creep resistance, and embrittlement behavior of steels and superalloys. In Fe--Mo alloys, the relevant structural landscape includes the comparatively compact A15, C15, C14, C36, $\sigma$, $\chi$, and $\mu$ phases together with the more intricate $R$, $M$, $P$, and $\delta$ phases. The physical challenge is not simply the multiplicity of competing prototypes, but the proliferation of distinct Wyckoff sublattices that permit chemically nontrivial order. Simple TCP structures already contain between two and five crystallographically distinct sites, whereas the complex phases spread the chemistry over 11--14 nonequivalent sublattices and unit cells as large as roughly 159 atoms. % TODO: CITE [general TCP review for transition-metal alloys] % TODO: CITE [Fe--Mo phase-diagram review] % TODO: CITE [experimental Fe--Mo TCP structures]

The resulting combinatorial space frustrates brute-force first-principles sampling. A chemically ordered prototype with many inequivalent sites admits a large number of distinguishable occupations even at fixed composition, and each candidate requires a full structural relaxation before its energetic ranking can be assessed. Although density-functional theory (DFT) is accurate enough to resolve the small energy differences that separate many TCP variants, the cost of exhaustive enumeration rises rapidly with unit-cell size and with the number of chemically distinct sites. This is precisely the regime in which data-efficient surrogate models are attractive, but it is also the regime in which naively flexible machine-learning representations are most vulnerable to overfitting because the available training data are necessarily sparse. % TODO: CITE [DFT studies of TCP phases in transition-metal alloys] % TODO: CITE [machine learning for intermetallics or crystalline materials with limited data]

The Fe--Mo repository analyzed here tackles this problem by embedding domain knowledge at three levels before regression is attempted. First, chemistry enters through Vegard-type volume scaling that generates sensible initial configurations across composition. Second, crystallography enters through coordination-number-resolved averaging (CNAV), which preserves the distinction between different local polyhedral environments instead of collapsing every structure into a single indiscriminate mean. Third, local bonding enters through bond-order-potential (BOP) moments, whose connection to the local electronic structure gives them a direct physical interpretation; ACE and SOAP descriptors serve as less specialized comparison baselines. The central question is then whether a training set drawn from the simpler TCP structures contains enough transferable information to predict the formation energies of the more complex phases and to support finite-temperature sublattice-occupancy analysis.

The archived notebooks indicate that the curated Fe--Mo dataset contains 292 usable structures after curation, that the downstream learning analysis scores 261 non-bcc/fcc/hcp structures on a phase-stratified 80/20 split, and that a separate complex-phase validation set contributes 70 additional DFT calculations distributed over $R$, $P$, $M$, and $\delta$. The internal test errors reach the low-\SI{10}{\milli\electronvolt\per\atom} range for the best domain-informed models, while the external parity plots for the complex phases remain in the few-times-\SI{10}{\milli\electronvolt\per\atom} regime. This combination of accuracy and transferability motivates a closer examination of which parts of the domain knowledge matter most and of how far one can push a modest DFT data budget.

In what follows, Sec.~\ref{sec:methods} describes the DFT data curation, the formation-energy definition, the descriptor construction, the stratified learning protocol, and the Bragg--Williams thermodynamic treatment. Section~\ref{sec:results} then presents the archived test-set errors, quantifies the effect of CNAV, summarizes the complex-phase validation workflow, and discusses the finite-temperature occupancy analysis for the $R$ phase. Section~\ref{sec:discussion} interprets the relative performance of the different descriptors and outlines the generality and limitations of the present approach before Sec.~\ref{sec:conclusions} closes with the implications for data-efficient modeling of chemically complex intermetallics.
